\chapter{Sampling Data}

\section{Planning a Study}

Every day, decisions get made based on data. A city council votes on a new transit route based on ridership surveys. 
A school principal adjusts the lunch schedule based on student feedback. A pharmaceutical company decides 
whether to release a new drug based on clinical trial results. In every case, the quality of the decision 
depends entirely on the quality of the data behind it.

To collect data well, you need a good plan! That plan must specify what question you are trying to answer, 
how you will gather information, and who you will gather this information from. 

% FIXME image/illustration needed

\subsection{Vocabulary}

Five terms form the foundation of everything that follows. Learn to distinguish them from each other, because 
you will probably be tested on this.

\begin{description}
    \item[\newterm{population}\index{population}] The entire group of individuals or items that 
    you want to learn about. The population is defined by the question you are asking.

    \item[\newterm{frame}\index{frame}] Also called a sampling frame. A list of all the members 
    of the population from which a sample will be drawn. The frame is usually close to the population, 
    but the two are not always identical. A school's enrollment database, a city's voter registration list, 
    and a company's customer records are all examples of frames.

    \item[\newterm{sample}\index{sample}] The part of the population that is actually studied. 
    You collect data from the sample and use it to draw conclusions about the population.

    \item[\newterm{sample survey}\index{sample survey}] The process of collecting information from a 
    sample. The information gathered is then used to make inferences about population parameters.

    \item[\newterm{census}\index{census}] The process of collecting information from every member of the population 
    rather than just a sample. A census gives complete information, but it comes at a cost.
\end{description}

The relationship between these terms is important. The population is the full group you care about. 
The frame is your working list of that group. The sample is the slice of the frame you actually study. 
A census is what happens when the sample and the population are the same thing or equal in number. 

% FIXME image/illustration needed

\begin{Exercise}[title={Identifying the Population and Sample}, label={ex:pop-sample-1}]
A city's parks department wants to understand how often residents use public parks. Staff 
randomly select 200 households from the city's water billing records and conduct phone interviews.

\begin{enumerate} 
    \item What is the population?
    \item What is the frame?
    \item What is the sample?
    \item Is this a census? Explain.
\end{enumerate}
\end{Exercise}
\begin{Answer}[ref={ex:pop-sample-1}]
\begin{enumerate} 
    \item The population is all households in the city.
    \item The frame is the city's water billing records.
    \item The sample is the 200 households selected and interviewed.
    \item No. A census would require interviewing every household in the city. Only 200 were selected here.
\end{enumerate}
\end{Answer}

\subsection{Why Not Always Use a Census?}

A census tells you about everyone in the population. That sounds great, but it is not always realistic. If the group is small, 
like a football team, you can measure everyone.

% FIXME image/illustration needed

Most of the time, a census is too difficult because:

- The population is too large. A company that wants to know what all households watch cannot contact every single one.
- The measurement destroys the item. If testing a battery kills it, checking every battery would waste them all.
- The information gets old too fast. If a company waits months to ask everyone, the results may no longer match what people think.

\begin{Exercise}[title={Census or Sample?}, label={ex:census-or-sample}]
For each situation, decide whether a census is feasible or whether a sample should be used. Briefly explain your 
reasoning.

\medskip
\begin{tabular}{|p{0.58\textwidth}|p{0.32\textwidth}|}
\hline
\textbf{Situation} & \textbf{Census or Sample? Explain} \\
\hline
A school librarian wants to know which of her 8 student volunteers prefers morning or afternoon shifts. & \\[2.5em]
\hline
A battery manufacturer wants to test the lifespan of batteries by running them until they die. & \\[2.5em]
\hline
A government agency wants to estimate the average commute time of workers across a country of 40 million employed adults. & \\[2.5em]
\hline
A restaurant owner wants to know the dietary restrictions of the 22 guests attending a private event. & \\[2.5em]
\hline
\end{tabular}
\end{Exercise}
\begin{Answer}[ref={ex:census-or-sample}]
\begin{itemize}
    \item \textbf{Library volunteers:} Census. With only 8 volunteers, asking all of them takes almost no effort.
    \item \textbf{Battery lifespan:} Sample. Testing can destroy the batteries, so a census would leave nothing to sell.
    \item \textbf{Commute times:} Sample. The population of 40 million is far too large for a full census to be practical.
    \item \textbf{Dietary restrictions:} Census. With only 22 guests, collecting complete information is straightforward 
    and important for planning.
\end{itemize}
\end{Answer}

In most studies, sampling is the only practical option, and the quality of the conclusions 
you draw depends entirely on how well the sample was chosen. 

\section{Planning and Conducting Surveys}

Not all samples are fair. How you pick people can make the group represent the population well, or it can leave out important parts.

% FIXME image/illustration needed

\subsection{Biased Sampling}

A sample is \newterm{biased}\index{bias} if it favors some people or answers over others. 
Bias does not go away just because the sample is bigger. 

% FIXME image/illustration needed

These three common methods usually make biased samples.

\newterm{Judgmental sampling}\index{judgmental sampling} is when one person decides who is included based on their own idea 
of what looks representative. A school inspector who visits only classrooms she already thinks are good is using judgmental 
sampling. The choice is based on their opinion.

A \newterm{convenience sample}\index{convenience sample} uses whoever or whatever is easiest to reach. 
A journalist who interviews only people at the bus stop outside his office is using a convenience sample. Those people 
might not represent all commuters.

A \newterm{voluntary response sample}\index{voluntary response sample} happens when people choose whether to answer. 
An online poll on a news site will attract mostly visitors with strong opinions. People who do not care much usually 
do not respond.

All three skip chance. Sampling methods based on random selection give everyone a fair shot and usually make more 
trustworthy results.

\subsection{Simple Random Sampling}

A \newterm{simple random sample}\index{simple random sample} (SRS) is a sample chosen so every group of 
the same size has the same chance of being picked. 

There are two ways to do it.

With \newterm{sampling with replacement}\index{sampling!with replacement}, you put someone back in the pool 
after choosing them. They could be picked again.

With \newterm{sampling without replacement}\index{sampling!without replacement}, you do not put them back. 
Each person can only be chosen once.

To get an SRS, you need a real random method.

% FIXME image/illustration needed

\begin{itemize}
    \item Assign every member of the population a number. Use a random number table, entering at any random point and reading 
    numbers sequentially. Select members whose numbers appear, skipping values outside your range and any repeats, until you have 
    the sample size you need.
    \item Use a calculator. On a TI-84, \texttt{MATH} $\to$ \texttt{PRB} $\to$ \texttt{5:randInt(} generates random integers. 
    The command \texttt{randInt(1,500,30)} produces 30 random integers between 1 and 500. Ignore repeats.
    \item Use a physical randomization device. Write each member's identifier on an identical slip of paper, mix thoroughly in a 
    container, and draw without looking until the sample is complete.
\end{itemize}

The essential requirement in every case is that chance dictates who is selected, rather than convenience or judgment. 
This is what makes the sample fair and the results trustworthy.

\begin{Exercise}[title={Selecting an SRS}, label={ex:srs}]
A regional transit authority wants to survey 25 of its 300 bus drivers about workplace safety conditions.

\begin{enumerate} 
    \item Describe a complete procedure for selecting an SRS of 25 drivers using a random number table.
    \item A supervisor suggests selecting the 25 drivers who are easiest to reach during the next shift change. What type of 
    sample is this, and what bias might it introduce?
    \item A colleague suggests selecting 5 drivers at random from each of the authority's 5 depots. Is this an SRS? Explain.
\end{enumerate}
\end{Exercise}
\begin{Answer}[ref={ex:srs}]
\begin{enumerate} 
    \item Number the 300 drivers from 001 to 300. Enter a random number table at any point and read three-digit numbers. 
    Select any driver whose number falls between 001 and 300, skipping values from 301 to 999 and any numbers that have already been 
    selected, until 25 drivers have been chosen. Alternatively, use \texttt{randInt(1,300,25)} on a TI-84, ignoring repeats.
    \item This is a convenience sample. Drivers who are easy to reach at shift changes may work particular routes or shifts that 
    differ systematically from others. For example, they may have different safety exposure than drivers on overnight or rural routes.
    \item No. An SRS requires that every possible group of 25 drivers from all 300 has an equal chance of being chosen. This procedure 
    cannot produce a sample where all 25 come from the same depot, so not all groups of 25 are possible. This is a stratified sample, 
    not an SRS.
\end{enumerate}
\end{Answer}

\subsection{Other Methods of Random Sampling}

Simple random sampling is not always the most practical or most efficient approach. Three other probabilistic methods appear 
regularly on the AP exam.

In \newterm{systematic sampling}\index{systematic sampling}, you list the population, choose a random start, and then pick 
every $k$th person. 
For example, to sample 40 students from 800, start at a random number between 1 and 20, then take every 20th student. This 
is simple when you have a full list.

% FIXME image/illustration needed

In \newterm{stratified random sampling}\index{stratified random sampling}, the population is split into groups called 
strata. 
Each group should be similar inside, like students in the same grade or patients in the same ward. Then you take a 
random sample from each group. 
This makes sure every part of the population is represented. This method of random sampling is especially useful in medical 
studies, market research, 
political polling, educational testing, etc, where you want to make sure you have enough people from each group to compare them.

% FIXME image/illustration needed

In \newterm{cluster sampling}\index{cluster sampling}, the population is split into groups called clusters. 
Each cluster should be like a mini-version of the whole population. Then you randomly pick some clusters and study everyone inside 
them. 
This works well when people are spread out over a wide area and it may be difficult or expensive to reach everyone. For example, 
a company that wants to survey its customers across the country might randomly select 10 cities and then survey all customers 
in those cities.
% FIXME image/illustration needed


\textbf{It is important to understand the difference between strata and clusters.} 
\newterm{Strata}\index{strata} are groups that are different from each other but similar within, and you sample from all of them. 
\newterm{Clusters}\index{clusters} are groups that are similar to the whole population, and you sample only some of them but take 
everyone inside those.


In \newterm{quota sampling}, a non-probability sampling method where researchers select participants until predefined quotas 
for certain groups 
(ie: age, gender, ethnicity) are filled. The sample is structured to reflect specific population characteristics, but participants 
are not chosen randomly. 
For example, if a population is 60\% female and 40\% male, a researcher surveys people until they reach those proportions.

% FIXME image/illustration needed

In \newterm{convenience sampling}, a non-probability sampling method where participants are selected because they are easy to 
access or available. 
For example, surveying students who happen to be in a campus cafeteria as you are eating lunch.

Bonus method:
\newterm{Snowball sampling}\index{snowball sampling} is a non-probabilistic method used when the population of interest is 
difficult or impossible to reach through conventional sampling. The researcher identifies a small number of initial participants 
who meet the study criteria, and then asks those participants to refer others they know who also qualify. Each wave of referrals 
generates the next, like a snowball rolling downhill and growing as it goes. It is not a commonly tested method on the AP exam, 
but it is worth knowing as an example of a non-probabilistic sampling technique that can be useful in certain contexts, 
such as studying hidden or hard-to-reach populations (e.g., people experiencing homelessness, members of stigmatized groups, or 
individuals with rare diseases).

\begin{Exercise}[title={Identifying Sampling Methods}, label={ex:samplingid}]
Identify the sampling method used in each scenario: SRS, systematic, stratified, or cluster.

\medskip
\begin{tabular}{|p{0.60\textwidth}|p{0.30\textwidth}|}
\hline
\textbf{Scenario} & \textbf{Method} \\
\hline
A supermarket chain selects 10 of its 80 stores at random and surveys every employee in those stores. & \\[2em]
\hline
A university randomly selects 50 students from each of its four faculties to participate in a wellbeing survey. & \\[2em]
\hline
A water utility selects a random starting meter on its numbered list and then tests every 25th meter for contamination. & \\[2em]
\hline
A film festival puts all 900 registered attendee names in a database and uses software to randomly select 90. & \\[2em]
\hline
\end{tabular}
\end{Exercise}
\begin{Answer}[ref={ex:samplingid}]
\begin{itemize}
    \item \textbf{Supermarket:} Cluster sampling. The stores are clusters, which means entire clusters are included once selected.
    \item \textbf{University:} Stratified random sampling. The four faculties are strata. An SRS is taken from each.
    \item \textbf{Water utility:} Systematic sampling. Random start followed by every 25th unit.
    \item \textbf{Film festival:} Simple random sample. Every possible group of 90 attendees is equally likely.
\end{itemize}
\end{Answer}

\subsection{Bias in Surveys}

Even a good random sample can still give bad results if the survey is set up badly. Bias can happen at many steps, further down 
the line from when 
you pick people.

\newterm{Sampling error}\index{sampling error} is the normal difference between a sample and the population. 
Different samples can give different answers just by chance. This is not the same as bias. Sampling error gets smaller with a 
bigger sample.

The following types of bias are more serious because they do not shrink with larger samples. They push results systematically in a 
particular direction regardless of sample size.

% FIXME image/illustration needed

\newterm{Response bias}\index{response bias} occurs when respondents give answers that do not reflect their true beliefs or 
behaviors. 
This often happens when the topic is sensitive or when the presence of an interviewer creates social pressure. A survey on 
hand-washing 
habits conducted by a hygiene inspector at a food preparation facility will likely overstate how often workers wash their hands, 
because workers know what answer the inspector wants to hear.

\newterm{Nonresponse bias}\index{nonresponse bias} occurs when individuals selected for the sample cannot be reached or choose not 
to participate, and those non-respondents differ systematically from those who do respond. A survey mailed to parents about school 
satisfaction will be completed at higher rates by engaged parents than by disengaged ones. The result overstates satisfaction among 
the parent community as a whole.

\newterm{Undercoverage bias}\index{undercoverage} occurs when part of the population is excluded from the sampling frame entirely. 
A survey about internet access conducted through online banner advertisements cannot reach people who lack internet access. Those 
excluded are precisely the group most relevant to the question.

\newterm{Wording effect bias}\index{wording effect bias} occurs when the phrasing of a question influences the answer. 
Consider the difference between these two ways of asking essentially the same thing:

\begin{itemize}
    \item  "Do you support stricter emissions standards for vehicles?"
    \item  "Do you support stricter emissions standards for vehicles, even if it means higher car prices and fewer model choices?"
\end{itemize}

The second version frames the question with costs attached, making a  "no" response far more likely even among people who 
genuinely support the policy. Neutral, direct wording is essential for unbiased responses.

\begin{Exercise}[title={Diagnosing Bias}, label={ex:bias-types}]
Identify the most likely type of bias in each scenario and explain why it applies.

\begin{enumerate} 
    \item A streaming platform sends a satisfaction survey to all subscribers but receives responses from only 8\% of them. 
    The platform reports that 91\% of subscribers are satisfied.

    \item A survey on recycling habits is conducted by going door-to-door in affluent neighborhoods only, because those streets 
    are easier and safer to walk.

    \item An interviewer asks employees:  "Our new open-plan office is designed to encourage collaboration. How much has it 
    improved your productivity?"

    \item A survey on alcohol consumption is conducted at a community health fair by a team wearing hospital uniforms. 
    Respondents significantly underreport their weekly intake.
\end{enumerate}
\end{Exercise}
\begin{Answer}[ref={ex:bias-types}]
\begin{enumerate} 
    \item Nonresponse bias. The 92\% who did not respond are unlikely to be satisfied in the same proportions as those who did. Highly satisfied and highly dissatisfied subscribers are both more motivated to respond than those with neutral feelings, but the pattern of non-response is unknown and the reported figure cannot be trusted.
    \item Undercoverage bias. Residents of lower-income neighborhoods are excluded from the frame entirely. Recycling behavior may differ systematically by neighborhood type, so the results cannot generalize to the full population.
    \item Wording effect bias. The question assumes the office has improved productivity and asks only  "how much." Respondents are not given a neutral opportunity to say productivity has not changed or has declined.
    \item Response bias. The formal appearance of the interviewers creates social pressure to give a health-conscious answer. Respondents adjust their reported behavior to match what they believe the interviewers expect.
\end{enumerate}
\end{Answer}



